\chapter{Especificaciones Técnicas}

\section{Tipo y categoría de cables}

\section{Normas de instalación, etiquetado y terminación}

\section{Estándares de racks y gabinetes}

\section{Equipos pasivos y activos}
La infraestructura del sistema de cableado estructurado requiere componentes
activos y pasivos que permitan garantizar conectividad, organización y adecuada
administración de la red.
% -----------------------------------------------------------
\subsection{Equipos pasivos}
% -----------------------------------------------------------
Los equipos pasivos corresponden a los elementos físicos del sistema de
telecomunicaciones que no requieren alimentación eléctrica para su funcionamiento.
% -----------------------------------------------------------
\subsubsection{Cable UTP Categoría 6}
% -----------------------------------------------------------

Se utilizará cable UTP Categoría 6 para el cableado horizontal del edificio,
permitiendo velocidades de transmisión de hasta 1 Gbps y frecuencias de
operación de 250 MHz.

\begin{table}[H]
\centering
\caption{Especificaciones del cable UTP Cat.6}
\begin{tabular}{|c|c|}
\hline
\textbf{Característica} & \textbf{Especificación} \\
\hline
Tipo & UTP Categoría 6 \\
\hline
Ancho de banda & 250 MHz \\
\hline
Velocidad soportada & 1 Gbps \\
\hline
Norma & ANSI/TIA-568-C.2 \\
\hline
\end{tabular}
\end{table}

% -----------------------------------------------------------
\subsubsection{Patch Panels}
% -----------------------------------------------------------

Los patch panels permiten organizar y administrar las terminaciones del cableado
horizontal dentro del rack principal de telecomunicaciones.

\begin{table}[H]
\centering
\caption{Características de los Patch Panels}
\begin{tabular}{|c|c|}
\hline
\textbf{Característica} & \textbf{Especificación} \\
\hline
Tipo & Patch Panel Cat.6 \\
\hline
Capacidad & 24 puertos \\
\hline
Montaje & Rack 19'' \\
\hline
Cantidad proyectada & 3 unidades \\
\hline
\end{tabular}
\end{table}

% -----------------------------------------------------------
\subsubsection{Faceplates y Keystone Jack}
% -----------------------------------------------------------

Los puntos de usuario estarán implementados mediante faceplates dobles con
conectores RJ-45 categoría 6, permitiendo conexión para equipos de datos y voz.

\begin{table}[H]
\centering
\caption{Características de faceplates y keystone jack}
\begin{tabular}{|c|c|}
\hline
\textbf{Elemento} & \textbf{Especificación} \\
\hline
Faceplate & Doble puerto RJ-45 \\
\hline
Keystone Jack & Cat.6 RJ-45 \\
\hline
Compatibilidad & ANSI/TIA-568-C \\
\hline
\end{tabular}
\end{table}

% -----------------------------------------------------------
\subsubsection{Rack de telecomunicaciones}
% -----------------------------------------------------------

El MDF contará con un gabinete metálico de 42U para alojar equipos activos y
pasivos de telecomunicaciones.

\begin{table}[H]
\centering
\caption{Características del rack principal}
\begin{tabular}{|c|c|}
\hline
\textbf{Característica} & \textbf{Especificación} \\
\hline
Tipo & Rack cerrado 19'' \\
\hline
Altura & 42U \\
\hline
Ancho & 600 mm \\
\hline
Profundidad & 800 mm \\
\hline
Ventilación & Ventiladores internos \\
\hline
\end{tabular}
\end{table}

% -----------------------------------------------------------
\subsection{Equipos activos}
% -----------------------------------------------------------

Los equipos activos son dispositivos electrónicos encargados de procesar,
administrar y distribuir el tráfico de red dentro de la infraestructura de
telecomunicaciones.

% -----------------------------------------------------------
\subsubsection{Switch de distribución}
% -----------------------------------------------------------

El MDF contará con un switch administrable Gigabit Ethernet encargado de la
distribución principal de la red del edificio.

\begin{table}[H]
\centering
\caption{Características del switch de distribución}
\begin{tabular}{|c|c|}
\hline
\textbf{Característica} & \textbf{Especificación} \\
\hline
Puertos & 24 puertos Gigabit Ethernet \\
\hline
Uplinks & 4 puertos SFP+ \\
\hline
Administración & Layer 2 administrable \\
\hline
Alimentación PoE & Compatible PoE+ \\
\hline
Velocidad & 10/100/1000 Mbps \\
\hline
\end{tabular}
\end{table}

% -----------------------------------------------------------
\subsubsection{Access Points Wi-Fi}
% -----------------------------------------------------------

Se instalarán Access Points Wi-Fi 6 para brindar cobertura inalámbrica en áreas
académicas y administrativas del edificio.

\begin{table}[H]
\centering
\caption{Características de los Access Points}
\begin{tabular}{|c|c|}
\hline
\textbf{Característica} & \textbf{Especificación} \\
\hline
Estándar & IEEE 802.11ax \\
\hline
Bandas & 2.4 GHz y 5 GHz \\
\hline
Tecnología & MU-MIMO y OFDMA \\
\hline
Alimentación & PoE+ \\
\hline
Capacidad & Hasta 500 usuarios \\
\hline
\end{tabular}
\end{table}

% -----------------------------------------------------------
\subsubsection{UPS}
% -----------------------------------------------------------

El rack principal contará con un sistema UPS encargado de proteger los equipos
de telecomunicaciones frente a cortes eléctricos y variaciones de tensión.

\begin{table}[H]
\centering
\caption{Características del UPS}
\begin{tabular}{|c|c|}
\hline
\textbf{Característica} & \textbf{Especificación} \\
\hline
Capacidad & 1500 VA \\
\hline
Autonomía & 30 minutos \\
\hline
Función & Respaldo eléctrico \\
\hline
Tipo & UPS interactiva \\
\hline
\end{tabular}
\end{table}
\section{Requisitos eléctricos y de puesta a tierra}

\section{Procedimientos de certificación}